\documentclass[11pt,twoside,openany]{book}

% This is a successor-side companion to the reconstructed 1995 manual.
% It deliberately does not modify or re-paginate manual/manual.tex.
\usepackage[letterpaper,margin=0.86in]{geometry}
\usepackage[T1]{fontenc}
\usepackage{lmodern}
\usepackage{microtype}
\usepackage{amsmath,amssymb}
\usepackage{array,booktabs,longtable}
\usepackage{enumitem}
\usepackage{xcolor}
\usepackage{listings}
\usepackage{fancyhdr}
\usepackage[hidelinks]{hyperref}

\definecolor{navy}{HTML}{17365D}
\definecolor{bluegray}{HTML}{475569}
\definecolor{paleblue}{HTML}{EEF5FB}
\definecolor{paleyellow}{HTML}{FFF8E7}
\definecolor{codegray}{HTML}{F5F7FA}

\hypersetup{
  pdftitle={TAORYX Language and Mathematics Reference},
  pdfauthor={TAORYX project},
  pdfsubject={Successor mathematics, problem grammar, and table syntax}
}

\setlist{nosep,leftmargin=*}
\setlength{\parindent}{0pt}
\setlength{\parskip}{5pt}
\renewcommand{\arraystretch}{1.18}
\newcommand{\taoryx}{TAORYX}
\newcommand{\pathcode}[1]{\path{#1}}
\newcommand{\claimbox}[1]{%
  \par\noindent\colorbox{paleblue}{\parbox{0.94\linewidth}{#1}}\par
}
\newcommand{\warningbox}[1]{%
  \par\noindent\colorbox{paleyellow}{\parbox{0.94\linewidth}{#1}}\par
}

\lstdefinestyle{taoryxcode}{
  basicstyle=\ttfamily\small,
  backgroundcolor=\color{codegray},
  frame=single,
  rulecolor=\color{bluegray!35},
  columns=fullflexible,
  keepspaces=true,
  breaklines=true,
  showstringspaces=false,
  tabsize=2
}
\lstset{style=taoryxcode}

\pagestyle{fancy}
\fancyhf{}
\fancyhead[L]{\small\textsc{TAORYX language and mathematics reference}}
\fancyhead[R]{\small\thepage}
\renewcommand{\headrulewidth}{0.4pt}

\title{\textcolor{navy}{TAORYX Language and Mathematics Reference}\\
  \large A successor-side companion to the TAOS User's Manual}
\author{TAORYX project}
\date{Successor documentation --- July 2026}

\begin{document}

\frontmatter
\maketitle
\thispagestyle{empty}

\claimbox{\textbf{Scope and relationship.} This document is organized to run
parallel to the reconstructed TAOS User's Manual: mathematical conventions
parallel the manual's Methods chapter, successor table syntax parallels Table
Files, and successor problem syntax parallels Problem Files. The original
manual remains authoritative for historical TAOS 96.0 wording and syntax. This
document records explicitly labelled \texttt{taoryx} successor behavior and
does not alter \pathcode{manual/manual.tex}.}

The companion documents have distinct jobs:

\begin{itemize}
  \item \pathcode{TAOS_manual_1995.pdf} and \pathcode{manual/} preserve the
    historical manual and its source traceability;
  \item \pathcode{docs/latex/taoryx_extensions_and_verification.tex} explains
    plant contracts, fidelity, controllers, scoring, and verification;
  \item this reference defines the successor language surface, equations,
    typed AST boundary, and example-run contract;
  \item \pathcode{docs/extensions/} and \pathcode{examples/taoryx/} provide
    the operational Markdown pages and executable fixtures.
\end{itemize}

\tableofcontents

\mainmatter

\chapter{Mathematical conventions and successor models}
\label{chap:successor-math}

\section{Claim boundary and notation}

The equations in this chapter are the semantic and runtime contracts used by
the successor examples. They are not a claim that every equation is present in
the 1995 executable, nor that the local runtime reproduces the historical
table library. Every implementation must declare its frames, signs, units,
state order, and coefficient conventions before numerical comparison.

The common variables are position $\boldsymbol r$, velocity
$\boldsymbol v$, body angular rate $\boldsymbol\omega_B$, mass $m$, density
$\rho$, relative speed $V_{\mathrm{rel}}$, reference area $S$, span $b$,
chord $c$, and dynamic pressure

\begin{equation}
  q_\infty = \tfrac{1}{2}\rho V_{\mathrm{rel}}^2.
\end{equation}

\section{Fidelity modes}

The successor profile exposes three related model choices. A mode is a model
selection, not an automatic assertion that a particular adapter can execute it.

\begin{longtable}{p{0.22\linewidth}p{0.27\linewidth}p{0.40\linewidth}}
\toprule
\textbf{Mode} & \textbf{Problem syntax} & \textbf{Semantic content}\\
\midrule
\endhead
Point mass & \pathcode{*3dof} or \pathcode{*mode point-mass} &
  Translational position and velocity with prescribed or reduced attitude and
  control effects.\\
Kinematic bridge & \pathcode{*mode kinematic-6dof} &
  Point-mass translation with prescribed or scheduled attitude variables; it
  is not free rotational dynamics.\\
Rigid-body 6-DOF & \pathcode{*6dof}, \pathcode{*sixdof}, or
  \pathcode{*mode rigid-body-6dof} &
  Translation, attitude, body rates, forces, moments, mass properties, and
  released control effects.\\
\bottomrule
\end{longtable}

The aliases are normalized in the AST while their source spelling remains
losslessly available: \pathcode{*3dof} denotes \texttt{point-mass}, and
\pathcode{*6dof} and \pathcode{*sixdof} denote \texttt{rigid-body-6dof}.

\section{Force, moment, and translational equations}

For a coefficient source using body-axis force coefficients and rolling,
 pitching, and yawing moment coefficients, dimensionalization is

\begin{align}
  \boldsymbol F_{A,B} &= q_\infty S
    \begin{bmatrix} C_X & C_Y & C_Z \end{bmatrix}^{T},\\
  \boldsymbol M_{A,B} &= q_\infty S
    \begin{bmatrix} b C_l & c C_m & b C_n \end{bmatrix}^{T}.
\end{align}

If the source uses wind-axis coefficients, the adapter must transform the
force into body axes before adding gravity and propulsion. The moment reference
must also be transferred explicitly:

\begin{equation}
  \boldsymbol M_{CG} = \boldsymbol M_{\mathrm{ref}}
    - \boldsymbol r_{\mathrm{ref}\to CG}\times\boldsymbol F_A.
\end{equation}

The body translational equation is

\begin{equation}
  m\left(\dot{\boldsymbol v}_B
    +\boldsymbol\omega_B\times\boldsymbol v_B\right)
  =\boldsymbol F_{A,B}+\boldsymbol F_{T,B}+\boldsymbol F_{g,B}.
\end{equation}

The full rigid-body rotational contract is

\begin{equation}
  I\dot{\boldsymbol\omega}_B
  +\boldsymbol\omega_B\times(I\boldsymbol\omega_B)
  =\boldsymbol M_{A,B}+\boldsymbol M_{T,B}+\boldsymbol M_{\mathrm{control},B}.
\end{equation}

The inertia tensor, quaternion convention, gravity model, and derivative
definitions are source contracts. They must not be inferred from coefficient
names or silently replaced by diagonal or constant values.

\section{Mass transitions and deployment}

Variable-mass propagation declares its method and mass property updates in the
problem file. The minimum mass-flow relationship is

\begin{equation}
  \dot m = -\dot m_{\mathrm{prop}},
\end{equation}

with CG and inertia updates supplied by the active mass model. A successor
deployment can initialize a new trajectory from a prior trajectory and segment:

\begin{lstlisting}
*method rk45-var variable-mass
*trajectory 2 payload start on 1
  *deployed from trajectory 1, segment 2, wt=250
\end{lstlisting}

The deployment source and optional weight are typed AST fields. Runtime support
still depends on the selected adapter and available mass properties.

\chapter{Problem files and grammar}
\label{chap:successor-problem}

\section{Profiles and claim modes}

TAORYX has two explicit profiles. \texttt{taos96} is the default historical
profile; \texttt{taoryx} opts into the successor surface. The parser does not
silently reinterpret successor directives as TAOS syntax.

\begin{lstlisting}[language=Python]
from taoryx.language import GrammarProfile, parse_problem_file

document = parse_problem_file(
    "mission.prb", profile=GrammarProfile.TAORYX
)
\end{lstlisting}

The selected profile is recorded on \texttt{ProblemDocument.grammar\_profile}.
Parsing returns lossless source, typed AST nodes, and source-located
diagnostics. A parsed file is not automatically runtime-complete.

\section{Hierarchy and extension blocks}

The historical hierarchy remains the organizing structure:

\begin{center}
problem $\rightarrow$ trajectory $\rightarrow$ segment $\rightarrow$ block body.
\end{center}

Successor blocks add to that hierarchy without replacing the manual-bounded
blocks:

\begin{longtable}{p{0.31\linewidth}p{0.20\linewidth}p{0.36\linewidth}}
\toprule
\textbf{Syntax family} & \textbf{Scope} & \textbf{AST/runtime meaning}\\
\midrule
\endhead
\pathcode{*3dof}, \pathcode{*6dof}, \pathcode{*sixdof} & problem &
  \texttt{DofDirectiveBlock}; normalized dynamics mode.\\
\pathcode{*mode ...} & problem & \texttt{ModeBlock}; explicit mode selection.\\
\pathcode{*method ...} & problem & \texttt{ExtensionBlock}; method and options.\\
\pathcode{*cases} & problem & \texttt{ExtensionBlock}; case columns and rows.\\
\pathcode{*mass ...}, \pathcode{*ptmass} & segment &
  \texttt{ExtensionBlock}; mass-model transition.\\
\pathcode{*deployed ...} & trajectory &
  \texttt{ExtensionBlock}; deployment source and weight.\\
\pathcode{surface\_ref(...)} & define body & typed helper declaration.\\
\pathcode{surface\_azm(...)} and \pathcode{surface\_dist(...)} & expressions &
  expression AST helper calls.\\
\pathcode{*runtime ...} & problem & \texttt{RuntimeBlock}; shared runtime contract.\\
\bottomrule
\end{longtable}

The full syntax and recovery behavior are maintained in the focused Markdown
reference, the documentary EBNF, and the typed parser sources:

\begin{itemize}
  \item \pathcode{docs/extensions/taoryx-language-reference.md};
  \item \pathcode{grammars/taoryx\_extensions.ebnf}; and
  \item \pathcode{src/taoryx/language/}.
\end{itemize}

The documentary EBNF is intentionally a
compact extension catalog; parser behavior and typed contracts are the
executable source of truth.

\section{Representative problem file}

\begin{lstlisting}
(successor-example)
*sixdof
*method rk4-fixed variable-mass
*atmos rcc ktf-annual
*wind geodetic file=wind.dat units=ft/sec

*trajectory 1 vehicle start on 1
  *initial ecic x=20925646 y=0 z=0 xdt=0 ydt=300 zdt=0 time=0 mass=100
  *segment 1 powered
    *mass xcg=0 ixx=10 iyy=11 izz=12
    *when time>1 stop

*trajectory 2 payload start on 1
  *deployed from trajectory 1, segment 1, wt=25
  *segment 1 coast
    *ptmass
    *when alt<0 stop
*end
\end{lstlisting}

\section{Validation and execution}

Focused syntax and negative-diagnostic fixtures use the successor profile:

\begin{lstlisting}[language=bash]
taoryx-validate --profile taoryx \
  examples/taoryx/syntax_fragments/*.prb \
  examples/taoryx/syntax_fragments/*.tbl
python tools/dev.py test-grammar
\end{lstlisting}

The local runtime can execute the explicitly supported subset and produce
normalized artifacts, telemetry, and plots:

\begin{lstlisting}[language=bash]
PYTHONPATH=src .venv/bin/python examples/run_corpus.py \
  --family taoryx --execute --output artifacts/examples/taoryx
\end{lstlisting}

\chapter{Table files and successor syntax}
\label{chap:successor-tables}

\section{Relationship to historical table files}

The successor table families use the same source-preserving philosophy as the
historical table parser: table names, headers, axes, assignments, operations,
comments, and diagnostics remain inspectable. New families are opt-in
successor syntax and must not be added to the TAOS 96 claim by catalog drift.

\begin{longtable}{p{0.25\linewidth}p{0.25\linewidth}p{0.37\linewidth}}
\toprule
\textbf{Table family} & \textbf{Typical independent axis} & \textbf{Purpose}\\
\midrule
\endhead
\pathcode{aero\_force} & Mach, alpha, beta, or control &
  Force coefficient channels such as $C_X$, $C_Y$, and $C_Z$.\\
\pathcode{aero\_moment} & alpha, beta, rates, or control &
  Moment coefficient channels such as $C_l$, $C_m$, and $C_n$.\\
\pathcode{inertia} & mass, time, or configuration &
  Mass-dependent inertia or mass-property channels.\\
\pathcode{cmx}, \pathcode{cmy}, \pathcode{cmz} & source-defined &
  Reviewed moment-table type aliases retained by the successor contract.\\
\bottomrule
\end{longtable}

\newpage
\section{New table declaration form}

The compact successor form declares an identification name, a successor type,
an output channel, an independent axis, and optional reference parameters:

\begin{lstlisting}
(force-coefficients) table aero_force ca(mach) sref=5.241 lrefx=2.58333
start
add ca(mach)
mach=0.0 1.0 2.0
ca=0.20 0.18 0.16
end

(mass-properties) table inertia ixx(mass)
start
add ixx(mass)
mass=0.0 100.0 200.0
ixx=1.0 10.0 18.0
end
\end{lstlisting}

The table parser retains the original source while typing the declaration,
axis, assignment, and operation records. It checks table-family membership,
axis order, assignment cardinality, supported header parameters, duplicate
assignments, and references supplied by problem files.

\section{Table-to-plant contract}

Before a table is used by a successor runtime, the adapter must declare the
coefficient convention and dimensionalization:

\begin{enumerate}
  \item identify the table family and coefficient frame;
  \item identify units and reference dimensions such as $S$, $b$, and $c$;
  \item declare interpolation, clamping, and out-of-envelope behavior;
  \item preserve table name and source hash in the run artifact; and
  \item verify that every problem-file reference resolves to the expected
    table family.
\end{enumerate}

The table type alone does not establish a physical truth claim. Source
provenance, valid ranges, and adapter behavior remain separate verification
inputs.

\chapter{Examples, artifacts, and traceability}
\label{chap:successor-examples}

\section{Corpus layout}

The examples are intentionally parallel to the manual's fragments and full
examples:

\begin{itemize}
  \item \pathcode{examples/taoryx/syntax\_fragments/} contains one focused
    fixture per extension family;
  \item \pathcode{examples/taoryx/full\_examples/} combines several features
    into realistic problem files;
  \item \pathcode{examples/taoryx/manifest.yaml} records profile, claim
    boundary, source paths, and execution expectations; and
  \item \pathcode{examples/taoryx/output\_contract.yaml} defines the common
    report, artifact, telemetry, and plot products.
\end{itemize}

\section{Traceability map}

\begin{longtable}{p{0.29\linewidth}p{0.57\linewidth}}
\toprule
\textbf{Question} & \textbf{Authoritative location}\\
\midrule
What did the original manual say? & \pathcode{manual/} and the source PDF.\\
What successor syntax is accepted? & \pathcode{src/taoryx/language/}, the grammar contracts, and focused fixtures.\\
What syntax is documented? & \pathcode{docs/extensions/taoryx-language-reference.md} and this reference.\\
What ran? & \pathcode{run.json}, \pathcode{artifact.json}, and \pathcode{telemetry.csv}.\\
What can be plotted? & The plot directory under the declared output contract.\\
\bottomrule
\end{longtable}

\warningbox{\textbf{Compatibility wording.} The local runtime executes the
explicitly supported TAOS96-compatible and TAORYX successor subsets. The
repository does not contain the historical executable and complete table
library, so neither this reference nor the verification guide claims bit-for-
bit TAOS 96.0 behavioral compatibility.}

\backmatter
\chapter*{Build commands}

Build this manual-like language reference with:

\begin{lstlisting}[language=bash]
python tools/dev.py language-reference
\end{lstlisting}

The normalized PDF is written to
\pathcode{output/pdf/taoryx_language_reference.pdf}. The composite successor
PDF includes this reference, the verification guide, and the Markdown
extension appendix:

\begin{lstlisting}[language=bash]
python tools/dev.py taoryx-extension-pdf
\end{lstlisting}

\end{document}
